\section{Finite calibration of the joint spectral fibre}\label{sec:v101-finite}
The full-rank inverse imposes a restriction on the possible leading
fibres which is absent for general compact semialgebraic sets. On each
multiplicity stratum a finite list of scalar, constrained observation
costs determines the entire joint spectral fibre, not merely its
coordinate radii. The constraints specify root sums and variances;
mixture weights and both stochastic channels are minimized over.

Throughout this section the centre belongs to $\mathfrak M_d$ of
Section~\ref{sec:multiplicity}. Distinct clusters are marked by their
limiting root and component; their multiplicities and endpoint status
form the pattern $\mathcal P$. The marking is unambiguous near the
centre, up to the already identified component permutation. Fix a small
closed cluster neighbourhood $B_\theta$ containing the centre in its
relative interior. No stochastic coordinate is held fixed in this
neighbourhood. Lemma~\ref{lem:fullrank} implies that every sufficiently
small-cost competitor in the whole closed model belongs to this
neighbourhood, after the component permutation.

For the roots $r_j+u_{j,i}$ in cluster $j$ put
\[
 s_j=\sum_{i=1}^{m_j}u_{j,i},\qquad
 X_j=\frac12\sum_{i=1}^{m_j}(u_{j,i}-s_j/m_j)^2.
\]
Write $\epsilon_j=1$ at $0$ and $\epsilon_j=-1$ at $D$; thus
$\epsilon_j s_j\ge0$ at an endpoint. There are two retained coordinate
systems.

If an interior cluster is repeated, retain the $X_j$ at all repeated
interior clusters and the $\epsilon_j s_j$ at all endpoint clusters.
Denote their number by $k$ and their cone by $\mathcal K=\R_+^k$.
All interior sums and all stochastic increments are free variables.
If no interior cluster is repeated, retain all cluster sums, with the
same sign reversal at $D$. Then
$\mathcal K=\R^{k_{\mathrm{int}}}\times\R_+^{k_{\mathrm{end}}}$,
and only the stochastic increments are free. In both cases $k\le2d$.
Let $\gamma=1/2$ in the first case and $\gamma=1$ in the second.

For $v\in\mathcal K$ and sufficiently small $t>0$, let
$B_\theta(t;v)$ be the competitors in $B_\theta$ whose retained
coordinates are exactly $tv$. In particular, prescribing $X_j=t v_j$
does not prescribe the cluster mean or its individual root spacings.
Define the scalar calibration cost
\begin{equation}\label{eq:v101-calibration}
 c_\theta(v)=\lim_{t\downarrow0}\frac4{t^2}
       \min_{\xi\in B_\theta(t;v)}h_w(F_T(\xi),P_\theta)^2.
\end{equation}
Existence, independence of a sufficiently small $B_\theta$, and the
fact that these are actual closed-model alternatives are included in
the theorem below. The normalization $4$ agrees with
\eqref{eq:mult-expand}.

Insert the retained and free variables into the score
\eqref{eq:score}, using the polynomial increments
\eqref{eq:mult-jet} in the repeated regime and its first-sum terms
in the linear regime. Write the resulting score as $La+Dv$, with $a$ free. Set
\begin{equation}\label{eq:v101-profile-matrix}
 Q_\theta=D^{\mathsf T}WD-D^{\mathsf T}WL
       (L^{\mathsf T}WL)^{-1}L^{\mathsf T}WD,
 \qquad W=\diag(w_j/P_{\theta,je}).
\end{equation}
Thus the score is $La+Dv$. An empty free block contributes zero.
In the repeated regime this is exactly \eqref{eq:mult-Q}, with the
same signs and the same free stochastic directions.

\begin{theorem}[Finite scalar determination]\label{thm:v101-finite}
For every centre in $\mathfrak M_d$ the limit
\eqref{eq:v101-calibration} exists and
\begin{equation}\label{eq:v101-calibration-form}
 c_\theta(v)=\min_a\|La+Dv\|_{I,\theta}^2
            =v^{\mathsf T}Q_\theta v,\qquad Q_\theta>0.
\end{equation}
The convergence in \eqref{eq:v101-calibration} is uniform for $v$ in
compact subsets of $\mathcal K$ and for centres in compact subsets of
a fixed pattern stratum. Its error is $O(\sqrt t)$ in the repeated
regime and $O(t)$ in the linear regime.

The finite marked scalar datum
\begin{equation}\label{eq:v101-finite-datum}
 \mathfrak C_\theta=
 \bigl(\mathcal P;\ c_\theta(e_i),\
             c_\theta(e_i+e_j):1\le i<j\le k\bigr)
\end{equation}
contains $k(k+1)/2$ numbers and determines the complete leading
spectral fibre, with its bottleneck metric. It also determines its
joint leading cluster-polynomial coefficient image. All mixture
weights and channel parameters remain unknown in these calibrations.

More explicitly,
\begin{equation}\label{eq:v101-polarization}
 (Q_\theta)_{ii}=c_\theta(e_i),\qquad
 (Q_\theta)_{ij}=\frac{c_\theta(e_i+e_j)-c_\theta(e_i)
                                      -c_\theta(e_j)}2.
\end{equation}
Let $\mathcal V_\theta=\{v\in\mathcal K:v^{\mathsf T}Q_\theta v\le4\}$.
In the repeated regime decode $v$ into all ordered, zero-sum vectors
$y_j$ with $\frac12\sum_i y_{j,i}^2=v_{X_j}$ at the repeated interior
clusters, and zero vectors at every other cluster. Endpoint coordinates
of $v$ are existentially quantified. In the linear regime decode
interior sums into the corresponding single root increments; at an
endpoint use all ordered one-sided increments with signed sum $v_j$.
The union of these decoded vectors is precisely the leading fibre
$\mathcal L_\theta$ of the whole observation ball, scaled by $t^{-\gamma}$.
\end{theorem}

\begin{proof}
The coefficient inverse gives an injective derivative in monic
coefficient and stochastic coordinates. Within either retained/free
system the polynomial directions are independent: after division by
$f_a$, their partial fractions have distinct simple poles, and distinct
double poles exactly at the repeated interior clusters. Hence the
combined score matrix $(L,D)$ is injective. Orthogonal projection off
$\operatorname{ran}L$ shows that \eqref{eq:v101-profile-matrix} is
positive definite and proves the algebraic minimum in
\eqref{eq:v101-calibration-form}. This also bounds its minimizing free
variable $a$ uniformly on the indicated compact sets.

We next prove that this quadratic program is the limit of the actual
constrained minima, not a substitute experiment. For a repeated
interior cluster and a prescribed $X\ge0$, choose an ordered vector
$y$ with $\sum_i y_i=0$ and $\frac12\sum_i y_i^2=X$. Such a vector
exists for every $X$: for example take one entry $-(m-1)b$ and the
other $m-1$ entries $b$, where $b^2=2X/(m(m-1))$. Use the exact roots
\[
 r_j+t a_j/m_j+\sqrt t\,y_{j,i}.
\]
At a simple interior root use its order-$t$ sum. At an endpoint realize
the prescribed signed sum by equal one-sided displacements. Give the
weights and channels their prescribed order-$t$ stochastic increments.
Strict interior margins for these parameters and separation of the
clusters make all these alternatives feasible for small $t$. They
satisfy the retained constraints exactly. Their observation expansion is
\begin{equation}\label{eq:v101-calibrated-expansion}
 F_T(\xi)-P_\theta=t(La+Dv)+O(t^{1+\gamma}),\qquad
 \frac4{t^2}h_w(F_T(\xi),P_\theta)^2
       =\|La+Dv\|_{I,\theta}^2+O(t^\gamma),
\end{equation}
where in the linear case all displacements, including the individual
endpoint displacements, are $O(t)$. Taking the minimizing $a$ gives
the required upper bound.

The constrained sets are closed in the compact cluster neighbourhood
and are nonempty by this construction, so their minima are attained.
The upper bound makes every minimizing observation error $O(t)$.
Lemma~\ref{lem:fullrank} then makes all coefficient and stochastic
errors $O(t)$, uniformly on compact sets of centres and retained
vectors. Analytic coprime cluster factorization gives the same bound
on cluster coefficients. The first coefficient bounds every free
interior sum by $O(t)$. At an endpoint the displacements have one
sign, so their prescribed sum controls every displacement. At a
repeated interior root the retained centred variance bounds each
centred displacement by $O(\sqrt t)$. Thus every minimizer has bounded
free first-order variables and satisfies
\eqref{eq:v101-calibrated-expansion}, with exactly its prescribed $v$.
Its leading score norm is no smaller than the displayed quadratic
minimum. This proves the lower bound, the limit and the uniform errors.
The same argument forces every minimizer into every smaller cluster
neighbourhood for small $t$; the germ and limit do not depend on its
initial choice. No channel or weight was fixed during either bound.

Each $e_i$ and $e_i+e_j$ belongs to the relevant cone, including in the
linear regime. Polarization therefore recovers the whole matrix from
legal scalar calibrations. Theorem~\ref{thm:mult-atlas}, whose direct
root-expansion and recovery argument was used above, identifies the
whole-model leading sets by precisely the stated decoders and the
constraint $v^{\mathsf T}Q_\theta v\le4$. This proves sufficiency for
the joint set, not just its diameter. Conversely, the decoder explicitly
recovers every leading point and every pair of points. Replacing each
root block by $\prod_i(z-y_{j,i})$ gives its joint coefficient image;
continuity of the fixed-degree root map gives the inverse on ordered
real-rooted blocks. The bottleneck metric is the maximum coordinate
metric on these ordered blocks, since separated limiting clusters cannot
be matched across one another. This proves the metric assertion.
\end{proof}

\begin{remark}[The marking and the scope of invariance]\label{rem:v101-marking}
The datum is intrinsic to the observation germ with its fixed spectral
coordinate and marked multiplicity pattern. Reparametrizing free channel
or mean variables leaves $\operatorname{ran}L$ unchanged. Adding a free
score column to a retained column leaves its quotient class unchanged.
An isometry of the observation tangent metric transports all these
classes isometrically. These operations preserve the calibration datum
and its decoded metric fibre. A permutation of the marking permutes the
datum. An arbitrary change of target coordinate must transport the
spectral decoder and its metric; no invariance under an untransported
change of units is asserted. The datum need not be minimal: existential
elimination of endpoint coordinates may identify two different matrices
at the level of the spectral image. We assert sufficiency, not the false
converse that equality of spectral images always implies equality of
all the calibration numbers.
\end{remark}

\begin{corollary}[Quantitative finite reconstruction]\label{cor:v101-stability}
Use the same pattern and decoder for two positive definite matrices
$Q,Q'$, and suppose $Q\ge\lambda I$ and
$\|Q'-Q\|_{\mathrm{op}}\le\eta<\lambda$. Put
$\vartheta=\eta/\lambda$ and $\beta=\gamma/2$. Then
\begin{equation}\label{eq:v101-radial-stability}
 (1+\vartheta)^{-\beta}\mathcal L(Q)
 \subseteq\mathcal L(Q')
 \subseteq(1-\vartheta)^{-\beta}\mathcal L(Q).
\end{equation}
If $R=\max_{y\in\mathcal L(Q)}\|y\|_\infty$, their Hausdorff distance
is at most
\[
 R\max\{1-(1+\vartheta)^{-\beta},\
                  (1-\vartheta)^{-\beta}-1\}.
\]
Their diameters satisfy the same two multiplicative bounds as
\eqref{eq:v101-radial-stability}. If each number in
\eqref{eq:v101-finite-datum} is known with error at most $e$, the
matrix reconstructed by \eqref{eq:v101-polarization} has operator-norm
error at most $3ke/2$.
\end{corollary}
\begin{proof}
The hypotheses give $(1-\vartheta)Q\le Q'\le(1+\vartheta)Q$.
The corresponding cone sections therefore lie between the radial
rescalings by $(1+\vartheta)^{-1/2}$ and
$(1-\vartheta)^{-1/2}$. Their root decoders are homogeneous of degree
$\gamma$: variances scale as squared roots in the repeated regime,
and sums scale as roots in the linear regime. This proves
\eqref{eq:v101-radial-stability}. Compare each point with its radial
rescaling to obtain the Hausdorff estimate; diameter is monotone under
set inclusion and homogeneous under dilation. Finally a diagonal entry
has error at most $e$ and an off-diagonal entry at most $3e/2$.
The maximal absolute row sum bounds the operator norm by $3ke/2$.
\end{proof}

The preceding theorem is not the distance-function reconstruction of an
arbitrary compact set. It first derives a quadratic quotient from the
full binary coefficient inverse, then realizes a finite collection of
its probes by feasible polynomial-root experiments, and finally decodes
that quotient into a constrained joint spectrum. The metric
reconstruction Theorem~\ref{thm:reconstruction} remains useful
for leading sets outside this class; no quadratic shape is assumed in
the general real atlas.
